\section{Introduction}
% Motivate the area and synthesize the individual literature reviews.

Publishing interpretable models can support transparency \citep{rudin2019stop}, but
leaf class counts also reveal statistics about training data. DRAFT
\citep{ferry2024draft} uses these statistics to reconstruct random-forest training
sets with constraint programming. \citet{hsu2026rashomon} apply this attack to sparse
near-optimal trees from a Rashomon set, enumerated by TreeFARMS
\citep{xin2022treefarms}, and show that releasing more trees can lower reconstruction
error. Our project asks what this average error leaves out.

Our individual literature reviews identified complementary gaps. De Coster studied
uncertainty: most reviewed attacks return one dataset, while
\citet{ferry2024probabilistic} describe uncertainty for a single interpretable model.
\citet{dick2023confidence} rank records by their frequency across reconstructions
but leave the estimation of calibrated confidence for future work. Fuligni studied background knowledge:
DRAFT varies the number of known columns, while \citet{mehnaz2022sensitive} vary
which attributes are unknown in a different, black-box inference setting. This
motivates testing which known attributes help recover a fixed target from a
Rashomon release.

Both reviews also point to the need for informed baselines. Census reconstruction and attribute
inference studies show why an attack should be compared with what can already be
inferred from population data \citep{abowd2025census,jayaraman2022imputation}.
We combine the two review strands by measuring both the effect of known attributes
and the reliability of full-record claims. Code, data and results are available at
\url{https://github.com/francescofuligni/rtds-project}.
